\uuid{0XGr}
\titre{Étude de suite}
\auteur{bodin ; hanani ; mzari}
\organisation{exo7}
\licence{CC BY-NC-SA 4.0}
\datecreate{2020-01-20}
\niveau{L1}
\module{Analyse}
\chapitre{Suite}
\sousChapitre{Autre}
\difficulte{3}
\theme{Suites, étude de suite}
\exo7qcmid{lille-1.suites.026}

\contenu{
\begin{question}
\qclassification{121.00}
\qtags{motcle=étude de suite}
Soit $(u_n)$ une suite croissante. On suppose que $\displaystyle u_{n+1}-u_n\leq \frac{1}{2^n}$ pour tout $n\geq 0$. Que peut-on en déduire ?
\begin{answers}
    \bad{$(u_n)$ est divergente.}
    \good{$(u_n)$ est bornée et $u_0\leq u_n\leq u_0+2$.}
    \good{$(u_n)$ est convergente et $\displaystyle u_0\leq \lim _{n\to +\infty}u_n\leq u_0+2$.}
    \bad{$\displaystyle \lim _{n\to +\infty}u_n=+\infty$}
\end{answers}
\begin{explanations}
On vérifie par récurrence que
$$\displaystyle u_0\leq u_n\leq u_0+1+\frac{1}{2}+\frac{1}{2^2}+\dots +\frac{1}{2^{n-1}}=u_0+2-\frac{1}{2^{n-1}}.$$
Donc, $u_0\leq u_n\leq u_0+2$. Étant à la fois croissante est majorée, la suite $(u_n)$ est convergente et, par passage à la limite, $\displaystyle u_0\leq \lim _{n\to +\infty}u_n\leq u_0+2$.
\end{explanations}
\end{question}
}
